%% 
%% Copyright 2019-2024 Elsevier Ltd
%% Complex Article Service (CAS) LaTeX template
%% Formatted for: Biomedical Signal Processing and Control (BSPC)
%% Using the official els-cas-templates bundle
%% 
\documentclass[a4paper,fleqn]{cas-dc}

% BSPC uses numbered citations [1, 2]
\usepackage[numbers]{natbib}
\usepackage{amsmath}
\usepackage{amssymb}

\begin{document}
\let\WriteBookmarks\relax
\def\floatpagepagefraction{1}
\def\textpagefraction{.001}

% Running headers
\shorttitle{Multimodal Closed-Loop Emotional Control in 2D Space}
\shortauthors{M. Sorodoc}

% Article title
\title [mode = title]{Simulation-Based Multimodal Closed-Loop Emotional Control in 2D Valence-Arousal Space Using Wearable EEG, Autonomic Signals, and Multi-Actuator Model Predictive Control}

% Author and affiliation details
\author[1]{Matei Sorodoc}
\cormark[1]
\ead{matei.sorodoc@student.utcluj.ro}

\affiliation[1]{organization={Master's Program in Cyber-Physical Systems, Technical University of Cluj-Napoca (UTCN)},
                city={Cluj-Napoca},
                country={Romania}}

\cortext[cor1]{Corresponding author}

% Abstract (BSPC Guide for Authors: maximum 250 words)
\begin{abstract}
Closed-loop Affective Brain-Computer Interfaces (aBCIs) are a promising approach for mental health support, but most existing prototypes only regulate one-dimensional arousal or stress using a single audio channel. Relying on music alone introduces actuator coupling issues, such as the tempo-valence crossover where slowing down music to calm a user also lowers perceived valence. In this paper, we present a simulation-based framework for multimodal closed-loop emotion regulation in the continuous two-dimensional Valence-Arousal space. We implement a causal physiological signal processing pipeline that combines electroencephalography (EEG) with electrodermal activity (EDA) and blood volume pulse (BVP) photoplethysmography. Using recordings from $N = 32$ participants in the DEAP dataset ($1,280$ trials, $148,800$ valid control steps at $\Delta t = 0.5\text{ s}$), we compare the full $32$-channel BioSemi 10--20 layout against two sparse wearable subsets ($6$-electrode and $4$-electrode montages). All filters use causal transposed second-order sections with steady-state initialization and past-only sliding spectral windows, giving zero lookahead leakage ($\max|\Delta \mathbf{y}_k| = 0.0$) compared to $27.60$ units of leakage with acausal zero-phase filtering. Per-step execution time stays below $1.55\text{ ms}$ (mean $\le 0.91\text{ ms}$ across all montages), leaving over $99.6\%$ timing slack within the $500\text{ ms}$ control interval. This causal feature pipeline provides the empirical input for a multi-actuator controller that coordinates adaptive music, ambient lighting, and heartbeat haptics to decouple valence and arousal regulation.
\end{abstract}

% Keywords (BSPC Guide for Authors: 1 to 7 keywords separated by \sep)
\begin{keywords}
Affective Brain-Computer Interface \sep Closed-Loop Emotion Regulation \sep Multimodal Fusion \sep Model Predictive Control \sep EEG \sep Electrodermal Activity \sep Valence-Arousal Model
\end{keywords}

\maketitle

\section{Introduction}\label{sec:intro}

Emotions affect daily decision-making, attention, and mental well-being, and poor emotion regulation under workplace stress is closely linked to anxiety, depression, and burnout \cite{palfi2024adaptive}. Over the past two decades, Affective Computing and passive Affective Brain-Computer Interfaces (aBCIs) have made steady progress in decoding emotional states from brain and physiological signals \cite{wu2023affective}. Most of these systems use Russell's two-dimensional (2D) Circumplex Model of Affect, which maps emotions along two bipolar dimensions: valence (pleasure--displeasure) and arousal (calm--excited) \cite{russell1980circumplex}. Using signals such as electroencephalography (EEG), electrodermal activity (EDA), photoplethysmography (PPG), and heart rate variability (HRV), modern machine learning models can decode induced emotions above chance, though single-trial classification performance remains modest and sensitive to signal noise and inter-subject variability \cite{wu2023affective,koelstra2012deap,saganowski2022emognition}. Moreover, recent reviews of neurofeedback and affective BCI systems note that most studies still stop at open-loop recognition \cite{zhang2026closed}. Simply detecting that a user is stressed or upset without adjusting external stimulation leaves the regulatory loop incomplete \cite{ewing2016evaluation,lorenz2016automatic,sayal2025music}.

Closing this loop, where real-time physiological measurements continuously adjust external stimuli to guide the user toward a target state, shifts the focus from passive monitoring to active regulation \cite{ewing2016evaluation,lorenz2016automatic}. Several early prototypes have shown that closed-loop emotion or affect regulation is feasible, especially using adaptive music \cite{sayal2025music,ramirez2015musical,daly2016affective,ehrlich2019closed}. In an elderly pilot study ($N = 10$, with $n = 6$ pre/post clinical evaluations), Ramirez et al. \cite{ramirez2015musical} used prefrontal EEG beta/alpha ratios (across AF3, AF4, F3, and F4) and frontal alpha asymmetry (F4--F3) to modulate music loudness and tempo, observing reduced depression scores. Daly et al. \cite{daly2016affective} built an affective Brain-Computer Music Interface (aBCMI) combining EEG, ECG, skin conductance, and respiration to guide users ($N = 8$ online) along four target trajectories (\textit{Make Happy}, \textit{Calm}, \textit{De-stress}, and \textit{Excite}). Ehrlich et al. \cite{ehrlich2019closed} ($n = 5$ online) used a real-time EEG classifier to drive an algorithmic music generator and found causal bidirectional interaction with the musical feedback in right-frontal gamma activity (over FC6). More recently, Zhang et al. \cite{zhang2026mindmelody} proposed MindMelody (preprint), which links a deep-learning EEG decoder with a language model to adapt generated audio. On the control-engineering side, Khazaei et al. \cite{khazaei2024bayesian} developed Bayesian expectation-maximization decoders to infer hidden arousal and cognitive performance states from skin conductance and behavioral signals during music listening, while Fekri Azgomi and Faghih \cite{azgomi2022enhancement} showed that hidden stress states inferred from skin conductance impulses can be regulated in simulation (\textit{in silico}) using Fuzzy-Supervised Linear Quadratic Regulators (LQR) and Model Predictive Control (MPC).

Even with these encouraging results, three practical issues still prevent current closed-loop systems from achieving reliable 2D emotional control with wearable sensors:

\textbf{1. 1D control versus 2D multimodal dynamics:} Existing mathematical control frameworks (such as LQR and MPC in \cite{azgomi2022enhancement}) only regulate a 1D state (stress/arousal) estimated from a single peripheral modality (EDA). Likewise, the online closed-loop BCI implementation in \cite{ehrlich2019closed} mapped a scalar classifier output $s \in (0,1)$ identically to both musical axes ($\text{valence} = \text{arousal} = s$), constraining feedback along a single 1D diagonal rather than steering across all four quadrants of the 2D emotion space. Tracking both dimensions reliably requires multimodal fusion informed by benchmark datasets such as DEAP \cite{koelstra2012deap}, FACED \cite{chen2023large}, and Emognition \cite{saganowski2022emognition}. Prior studies show complementary roles across central and autonomic markers: frontal alpha asymmetry responds to valence shifts \cite{ramirez2015musical,koelstra2012deap}, midline frontal theta (Fz) and parietal upper alpha (P4) index mental effort and cognitive demand \cite{ewing2016evaluation}, EEG phase synchronization combined with ECG-derived HRV captures 2D affective quadrants \cite{marin2018affective}, and EDA impulse deconvolution supports continuous sympathetic arousal tracking \cite{khazaei2024bayesian}. Single-trial DEAP classification also confirmed complementary contributions between EEG and peripheral signals across arousal and valence \cite{koelstra2012deap}. However, film-induction studies show that affective arousal and physiological responses can transfer across consecutive stimuli (emotional carry-over) \cite{gross1995emotion,fernandez2019effective}, as quantified in sequential EEG datasets by the Affective Priming Score (APS) \cite{gutierrez2025affective}. Combining fast cortical EEG changes with slower autonomic EDA/HRV dynamics without modeling this sequential carry-over risks state-estimation errors right after emotional transitions.

\textbf{2. Coupling limitations of single-actuator music feedback:} Most closed-loop prototypes rely on a single feedback channel, usually adaptive music \cite{sayal2025music,daly2016affective,ehrlich2019closed}. However, recent validation studies of parametric music generators (AffectMachine-Classical \cite{agres2023affectmachine} and AffectMachine-Pop \cite{agres2025affectmachine}) revealed an important side effect: while music parameters control perceived arousal strongly ($R^2 = 0.86\text{--}0.96$) alongside high valence fit ($R^2 = 0.90\text{--}0.93$), perceived valence suffers from an unintended tempo-valence crossover. Specifically, when target valence is high but arousal is driven low (a calm, positive state), listeners perceive the slower music as less pleasant than intended, causing an unwanted drop in valence \cite{ehrlich2019closed,agres2023affectmachine,agres2025affectmachine}. Conversely, in Daly et al. \cite{daly2016affective}, single-modality monophonic piano feedback failed to raise arousal for most participants during \textit{Excite} tasks, often producing no change or decreased arousal alongside an unintended valence increase. Overcoming this actuator coupling motivates combining music with complementary non-acoustic modalities: (i) ambient colored lighting, where Li et al. \cite{li2025designing} ($N = 33$) showed that illuminance strongly modulates both perceived valence and arousal (with bright green light at 180 lx yielding the highest self-reported valence and dim purple light at 20 lx elevating EEG frontal midline theta and frontal alpha asymmetry markers associated with positive emotion regulation); and (ii) wrist-worn heartbeat haptics, where Azevedo et al. \cite{azevedo2017calming} ($N = 52$) showed that tactile pulses delivered 20\% slower than resting heart rate (clamped to 40--65 bpm) significantly reduced skin conductance arousal and state anxiety during stress anticipation, without changing mean heart rate. So far, no closed-loop control study has integrated music, lighting, and haptic feedback into a unified multi-actuator mathematical model.

\textbf{3. Lack of simulation benchmarks before human testing:} In biomedical control fields like automated anesthesia or insulin delivery, controllers are tested in computer simulations (\textit{in silico}) on virtual patient models before human deployment \cite{azgomi2022enhancement}. In affective computing, by contrast, many closed-loop prototypes are piloted directly on small volunteer cohorts ($N = 5\text{--}10$) \cite{zhang2026closed,sayal2025music,ramirez2015musical,daly2016affective,ehrlich2019closed}. General neurofeedback reviews report non-responder rates ranging from 16\% to 57\% (mean 38\%) \cite{sayal2025music}, while cross-session EEG non-stationarity and signal drift degrade decoder and controller stability across sessions \cite{zhang2026closed,wu2023affective}. A simulation benchmark built from real physiological datasets lets us compare model-based controllers (MPC and Fuzzy-Supervised MPC \cite{azgomi2022enhancement}) against model-free Reinforcement Learning (RL) \cite{pohlmeyer2014using} under realistic noise, temporary sensor dropouts, and ethical user-safety, agency, and mental-privacy constraints \cite{steinert2020wired}.

To address these gaps, this study asks the following central research question:
\begin{quote}
How can multimodal physiological signals from sparse wearable sensors be reliably fused to track emotional dynamics in real time, and how effectively can a coordinated multi-actuator controller steer a user state across the 2D Valence-Arousal space?
\end{quote}
We break this down into three sub-questions:
\begin{itemize}
    \item \textbf{RQ1 (Multimodal Observability):} Can combining sparse frontal and parietal EEG with peripheral EDA and HRV in a Kalman filter accurately track 2D emotional trajectories while attenuating slow autonomic lag and sequential carry-over?
    \item \textbf{RQ2 (Multi-Actuator Synergy):} Does coordinating parametric music with ambient colored lighting and slow wrist heartbeat vibrations mitigate the musical tempo-valence crossover, enabling independent steering of both valence and arousal?
    \item \textbf{RQ3 (Robustness):} How do optimal model-based controllers (MPC and Fuzzy-Supervised MPC) compare with data-driven Reinforcement Learning in tracking speed, control effort, and resilience against temporary EEG sensor dropouts?
\end{itemize}

To answer these questions, we present an \textit{in silico} multimodal closed-loop emotional control framework for the 2D valence-arousal space. Instead of using hand-tuned synthetic parameters, we identify data-driven 2D state-space Virtual Subjects from physiological recordings of $N = 32$ participants in the DEAP dataset \cite{koelstra2012deap}. To keep the setup realistic for portable wearable headsets, we benchmark the full 32-channel 10--20 EEG montage against two sparse wearable subsets: a 6-electrode configuration (Fz, F3, F4, P4, T7, T8) and a 4-electrode fronto-temporal configuration (AF3, AF4, T7, T8), each paired with peripheral EDA (GSR) and BVP/PPG-derived HRV signals.

The main contributions of this paper are:
\begin{itemize}
    \item \textbf{Multimodal 2D state estimator with carry-over filtering:} We design a discrete-time 2D Valence-Arousal state-space model and Kalman filter that combines fast EEG features with slow EDA/HRV features, incorporating Affective Priming Score (APS)-inspired \cite{gutierrez2025affective} sequential carry-over weighting to reduce transition errors.
    \item \textbf{Three-actuator coordination model:} We model the combined effect of parametric music \cite{ehrlich2019closed,agres2023affectmachine,agres2025affectmachine}, colored ambient lighting \cite{li2025designing}, and slow heartbeat wrist vibrations \cite{azevedo2017calming}, showing how lighting and haptics help compensate for the musical tempo-valence crossover.
    \item \textbf{Closed-loop controller comparison and dropout testing:} We compare Open-Loop stimulation, LQR, Constrained MPC, Fuzzy-Supervised MPC \cite{azgomi2022enhancement}, and model-free Reinforcement Learning \cite{pohlmeyer2014using} across four regulation tasks (\textit{De-stress}, \textit{Calm}, \textit{Make Happy}, and \textit{Excite} \cite{daly2016affective}), evaluating tracking accuracy, settling time, and robustness when EEG channels temporarily drop out due to movement.
\end{itemize}

The remainder of the paper is structured as follows. Section~\ref{sec:methods} describes the dataset, wearable montage selection, and causal multimodal feature extraction pipeline, while the following sections cover 2D state-space identification, multi-actuator modeling, and observer-controller design. Section 3 presents the \textit{in silico} simulation results for state tracking, actuator coordination, and sensor-dropout robustness. Section 4 discusses the findings, wearable hardware considerations, and ethical safeguards, and Section 5 concludes the paper.

\section{Materials and Methods}\label{sec:methods}

\subsection{Dataset and Wearable Montage Ablation Design}\label{subsec:dataset}

To ground state estimation in real human physiology rather than synthetic test signals, we used the DEAP multimodal database ($N = 32$ participants, 40 one-minute music-video trials per participant, totaling 1,280 trials at $f_s = 128\text{ Hz}$) \cite{koelstra2012deap}. Each preprocessed recording has tensor shape $(40, 40, 8064)$ spanning $63.0\text{ s}$: $3.0\text{ s}$ of pre-stimulus neutral fixation ($384$ samples, indices $0{:}384$) followed by $60.0\text{ s}$ of affective stimulation ($7,680$ samples, indices $384{:}8064$). In the preprocessed Python \texttt{.dat} distribution, all $439$ Quadrant-1 trials ($\text{Valence} > 5.0$, $\text{Arousal} > 5.0$) store inverted labels $(9.0 - \text{Valence}, 9.0 - \text{Arousal})$, mapping $40$ trials with $9.00$ ratings to $0.00$ ($31$ valence zeros and $15$ arousal zeros, with $6$ co-occurring) and $92$ trials to $(0.0, 1.0)$. Rather than discarding these trials, we restored the authoritative Self-Assessment Manikin (SAM) ratings directly from \texttt{participant\_ratings.xls} aligned by \texttt{Participant\_id} and \texttt{Experiment\_id} (cross-verified to $10^{-6}$ against \texttt{Dominance} and \texttt{Liking} in \texttt{.dat}), recovering all $1,280$ valid trials ($153,600$ valid control steps at $\Delta t = 0.5\text{ s}$) across $[1.00, 9.00]$ for valence (mean $5.25$, SD $2.13$) and arousal (mean $5.16$, SD $2.02$). Raw dataset files were kept read-only on disk and in memory, and verified against SHA-256 checksums before and after feature extraction.

Because 32-channel laboratory EEG caps are impractical for everyday wearable use, we compared the full $32$-channel BioSemi Geneva 10--20 layout against two sparse wearable montages. We used explicit channel indices rather than generic toolbox templates, which differ from DEAP's Geneva channel sequence at 30 of 32 positions. Specifically, we benchmarked: (i) full\_32 (channels $0\text{--}31$, $135$ EEG features: $32 \times 4$ spectral band powers plus $7$ bilateral alpha asymmetries across Fp2--Fp1, AF4--AF3, F4--F3, F8--F7, FC2--FC1, FC6--FC5, and T8--T7); (ii) wearable\_6 (indices $[18, 2, 19, 28, 7, 25]$ corresponding to Fz, F3, F4, P4, T7, and T8, $27$ EEG features: $6 \times 4$ band powers plus frontal alpha asymmetry [F4--F3], frontal theta-beta ratio [Fz], and parietal beta-alpha ratio [P4]); and (iii) wearable\_4 (indices $[1, 17, 7, 25]$ corresponding to AF3, AF4, T7, and T8, $18$ EEG features: $4 \times 4$ band powers plus AF4--AF3 and T8--T7 alpha asymmetries). Each EEG montage was combined with $11$ peripheral features extracted from finger Galvanic Skin Response (GSR, Hardware Channel 37 / Array Index 36) and thumb Blood Volume Pulse (BVP, Hardware Channel 39 / Array Index 38), giving observation dimensions $D = 146$ (full\_32), $D = 38$ (wearable\_6), and $D = 29$ (wearable\_4).

\subsection{Causal Real-Time Preprocessing and Multimodal Feature Extraction}\label{subsec:causal_features}

To avoid future-to-past data leakage (an acausal filtering ``inverse crime''), we did not use bidirectional zero-phase filters (filtfilt), whole-trial independent component analysis, or trial-wide $z$-scoring. Instead, every physiological stream passed forward sample-by-sample through Direct Form II Transposed Second-Order Sections (SOS) Infinite Impulse Response (IIR) Butterworth filters initialized at steady state ($\mathbf{z}_0 = \text{sosfilt\_zi} \cdot x[0]$) and propagated across consecutive $\Delta t = 0.5\text{ s}$ ($64$-sample) control steps with persistent state vectors $\mathbf{z}_k$. Filter states, baseline statistics, and sliding FIFOs were reset between trials. Before filtering, transient non-finite dropouts triggered causal sample-and-hold imputation, and impulse spikes were clamped to pre-stimulus baseline bounds ($\mu_{\text{base},c} \pm 5\sigma_{\text{base},c}$ when $\sigma_{\text{base},c} > 10^{-6}\ \mu\text{V}$) computed solely from samples $0{:}384$.

For EEG, a causal 4th-order Butterworth bandpass SOS filter ($4.0\text{--}45.0\text{ Hz}$, mean passband group delay $3.13$ samples [$24.42\text{ ms}$] at $f_s = 128\text{ Hz}$) conditioned selected channels before computing past-only $W_{\text{EEG}} = 2.0\text{ s}$ ($N_W = 256$ samples) Hamming-windowed rFFT periodograms over $[k - N_W + 1,\, k]$. Band powers $P_{k,c,b}$ [$\mu\text{V}^2$] were integrated over $\theta$ ($4\text{--}8\text{ Hz}$), $\alpha$ ($8\text{--}12\text{ Hz}$), $\beta$ ($12\text{--}30\text{ Hz}$), and $\gamma$ ($30\text{--}45\text{ Hz}$), log-transformed with a numerical floor $\epsilon = 10^{-12}\ \mu\text{V}^2$, and baseline-corrected against the $3.0\text{ s}$ fixation epoch via
\begin{equation}\label{eq:eeg_log_norm}
\Delta \tilde{P}_{k,c,b} = \ln(P_{k,c,b} + \epsilon) - \ln(P_{\text{base},c,b} + \epsilon)
\end{equation}
where $\Delta \tilde{P}_{k,c,b}$ [dimensionless] is the causal baseline-normalized log-bandpower at control step $k$, $P_{k,c,b}$ [$\mu\text{V}^2$] is the current window bandpower, and $P_{\text{base},c,b}$ [$\mu\text{V}^2$] is the pre-stimulus fixation bandpower.

For autonomic monitoring, we avoided subtractive phasic GSR decomposition ($y_{\text{raw}} - y_{\text{tonic}}$) because a 2nd-order $0.05\text{ Hz}$ lowpass filter has a mean group delay of $621.59$ samples ($4.86\text{ s}$ at $f_s = 128\text{ Hz}$), which creates phase-mismatch gain peaking when subtracted from undelayed raw GSR. Instead, tonic Skin Conductance Level (SCL) and phasic Skin Conductance Responses (SCR) were extracted separately using two parallel 2nd-order Butterworth SOS filters. A $0.05\text{ Hz}$ lowpass filter yields window-mean SCL [$\mu\text{S}$] and causal SCL slope [$\mu\text{S}/\text{s}$], while a $0.05\text{--}1.5\text{ Hz}$ bandpass filter (mean group delay $42.49$ samples [$331.98\text{ ms}$]) feeds a 1-sample-delayed causal peak detector with baseline-adaptive threshold $\Delta S_{\min} = 0.20\,\sigma_{\text{base,GSR}}$ and $1.0\text{ s}$ refractory lockout to compute SCR power [$\mu\text{S}^2$], peak rate [$\text{peaks/min}$], and mean SCR amplitude [$\mu\text{S}$] ($0.0$ fallback when zero peaks occur).

Thumb BVP was conditioned with a 2nd-order $0.5\text{--}5.0\text{ Hz}$ bandpass SOS filter (mean group delay $14.31$ samples [$111.78\text{ ms}$]), followed by a 1-sample-delayed systolic peak detector with rolling $2.0\text{ s}$ adaptive thresholding ($\gamma = 0.60$) and a $300\text{ ms}$ ($39$-sample) refractory gate to prevent heart rates above $200\text{ bpm}$. Across a rolling $W_{\text{HRV}} = 30.0\text{ s}$ buffer (expanding causally during cold start $t < 30.0\text{ s}$), we computed mean HR [$\text{bpm}$], SDNN [$\text{s}$], RMSSD [$\text{s}$], and Lomb--Scargle LF ($0.04\text{--}0.15\text{ Hz}$), HF ($0.15\text{--}0.40\text{ Hz}$), and LF/HF ratio whenever $M \ge 10$ beats and $t \ge 30.0\text{ s}$ elapsed. During cold start or peripheral vasoconstriction dropouts with $M < 2$ beats (which occurred in $4,070$ of $153,600$ cohort windows [$2.65\%$]), the BVP branch fell back deterministically to pre-stimulus baseline values and emitted a degraded quality confidence scalar $\omega_{\text{BVP}} = 0.1$ ($1.0$ otherwise), preventing any non-finite divisions.

\subsection{Deterministic Verification and Feature Storage}\label{subsec:phase1_verification}

Causal verification confirmed zero lookahead: injecting a $10^5\,\mu\text{V}$, $20\text{ Hz}$ burst at sample $1024$ produced $\max|\Delta \mathbf{y}_k| = 0.0$ across all causal streaming steps ending at or before $1024$, compared with $27.60$ log-power leakage units when the same pipeline used acausal sosfiltfilt. Under single-threaded BLAS execution, mean per-step processing time measured $0.903\text{ ms}$ ($\max = 1.547\text{ ms}$) for full\_32, $0.842\text{ ms}$ ($\max = 1.462\text{ ms}$) for wearable\_6, and $0.822\text{ ms}$ ($\max = 1.165\text{ ms}$) for wearable\_4, leaving over $99.6\%$ slack within the $500\text{ ms}$ control interval. All $18$ deterministic unit tests passed without error, covering SHA-256 dataset integrity, lookahead leakage, flatline, impulse, DC, and NaN stress stability, montage channel orthogonality, feature store staging, BVP cold start, and single-threaded bitwise reproducibility. The staged HDF5 and NPZ feature archives for all three montages ($153,600$ windows each) contained zero non-finite entries.

\printcredits

\bibliographystyle{cas-model2-names}
\bibliography{references}

\end{document}
